\documentclass[11pt]{article}

\usepackage[margin=1in]{geometry}
\usepackage[T1]{fontenc}
\usepackage{lmodern}
\usepackage{hyperref}
\usepackage{xcolor}
\usepackage{listings}
\usepackage{enumitem}
\usepackage{titlesec}

\hypersetup{
    colorlinks=true,
    linkcolor=blue,
    urlcolor=blue
}

\titleformat{\section}
{\large\bfseries}
{\thesection.}{0.5em}{}

\title{CPSC 250 Development Environment Setup}
\author{Christopher Newport University}
\date{Summer 2026}

\lstset{
    basicstyle=\ttfamily\small,
    breaklines=true,
    frame=single,
    backgroundcolor=\color{gray!10}
}

\begin{document}

\maketitle

\section{Purpose}

In CPSC 250, students will regularly run, modify, and experiment with Python code examples provided throughout the course. To support this workflow, all students should configure a working Python development environment on their personal laptop.

This document walks you through:

\begin{itemize}
    \item Installing Python
    \item Installing PyCharm
    \item Creating a GitHub account
    \item Forking the course repository
    \item Cloning the repository in PyCharm
    \item Creating a Python virtual environment
    \item Running example Python programs
    \item Using GitHub through PyCharm
\end{itemize}

\section{Required Software}

Students must install the following software:

\begin{itemize}
    \item Python 3
    \item PyCharm Community Edition
    \item Git
\end{itemize}

\section{Install Python}

\subsection{Windows}

\begin{enumerate}
    \item Go to:
    \[
    \texttt{https://www.python.org/downloads/}
    \]

    \item Download the latest version of Python 3.

    \item IMPORTANT: During installation, check:
    \begin{quote}
    ``Add Python to PATH''
    \end{quote}

    \item Complete the installation.
\end{enumerate}

\subsection{macOS}

\begin{enumerate}
    \item Go to:
    \[
    \texttt{https://www.python.org/downloads/}
    \]

    \item Download the latest macOS installer.

    \item Run the installer and complete installation.
\end{enumerate}

\section{Install PyCharm Community Edition}

\begin{enumerate}
    \item Go to:
    \[
    \texttt{https://www.jetbrains.com/pycharm/download/}
    \]

    \item Download:
    \begin{quote}
    PyCharm Community Edition
    \end{quote}

    \item Install using the default settings.
\end{enumerate}

\section{Install Git}

\subsection{Windows}

\begin{enumerate}
    \item Go to:
    \[
    \texttt{https://git-scm.com/download/win}
    \]

    \item Download and install Git using the default settings.
\end{enumerate}

\subsection{macOS}

Git is often already installed.

Open the Terminal application and type:

\begin{lstlisting}
git --version
\end{lstlisting}

If Git is not installed, macOS will usually prompt you to install the command line developer tools automatically.

\section{Create a GitHub Account}

\begin{enumerate}
    \item Go to:
    \[
    \texttt{https://github.com}
    \]

    \item Create a GitHub account if you do not already have one.

    \item Verify that you can successfully log in.
\end{enumerate}

\section{Fork the Course Repository}

The course repository contains lecture examples, sample programs, notes, and supporting materials used throughout the semester.

\subsection{Course Repository}

\[
\texttt{https://github.com/brash99/cpsc250}
\]

\subsection{Fork Instructions}

\begin{enumerate}
    \item Log into GitHub.

    \item Navigate to:
    \[
    \texttt{https://github.com/brash99/cpsc250}
    \]

    \item Click the ``Fork'' button near the upper-right corner.

    \item GitHub will create your own personal copy of the repository.
\end{enumerate}

\section{Clone Your Fork Using PyCharm}

\begin{enumerate}
    \item Open PyCharm.

    \item Select:
    \begin{quote}
    Get from VCS
    \end{quote}

    \item Log into GitHub when prompted.

    \item IMPORTANT:
    Select YOUR fork of the repository, NOT the instructor repository.

    \item Choose a local directory for the project.

    \item Click:
    \begin{quote}
    Clone
    \end{quote}
\end{enumerate}

\section{Create a Python Virtual Environment}

PyCharm will usually offer to create a virtual environment automatically.

Accept the default settings.

If needed, you can manually create one later.

\subsection{Windows PowerShell}

\begin{lstlisting}
python -m venv venv
\end{lstlisting}

Activate:

\begin{lstlisting}
.\venv\Scripts\Activate.ps1
\end{lstlisting}

\subsection{macOS Terminal}

\begin{lstlisting}
python3 -m venv venv
\end{lstlisting}

Activate:

\begin{lstlisting}
source venv/bin/activate
\end{lstlisting}

\section{Install Required Packages}

If the repository contains a \texttt{requirements.txt} file, install packages using:

\subsection{Windows PowerShell}

\begin{lstlisting}
pip install -r requirements.txt
\end{lstlisting}

\subsection{macOS Terminal}

\begin{lstlisting}
pip3 install -r requirements.txt
\end{lstlisting}

\section{Run an Example Program}

\begin{enumerate}
    \item Open one of the example Python files in the repository.

    \item Right-click inside the editor window.

    \item Select:
    \begin{quote}
    Run
    \end{quote}

    \item Verify that the program executes successfully.
\end{enumerate}

\section{Make a Small Test Change}

To verify that GitHub integration works correctly:

\begin{enumerate}
    \item Open a Python file.

    \item Add a simple comment such as:
\end{enumerate}

\begin{lstlisting}
# Testing my development environment
\end{lstlisting}

\section{Commit and Push Using PyCharm}

\begin{enumerate}
    \item In PyCharm, select:
    \begin{quote}
    Git $\rightarrow$ Commit
    \end{quote}

    \item Enter a commit message such as:
\end{enumerate}

\begin{lstlisting}
Test commit
\end{lstlisting}

\begin{enumerate}[resume]
    \item Commit the changes.

    \item Then select:
    \begin{quote}
    Git $\rightarrow$ Push
    \end{quote}

    \item Verify that the changes appear in your GitHub repository.
\end{enumerate}

\section{Windows PowerShell vs Git Bash}

These instructions use Windows PowerShell as the default Windows terminal environment because:

\begin{itemize}
    \item PowerShell is installed by default on modern Windows systems.
    \item It integrates well with PyCharm.
    \item It works well with GitHub and GitHub CLI.
    \item It minimizes additional installation complexity.
\end{itemize}

Some other courses may use Git Bash instead. Either environment can work successfully for Python development.

\section{Troubleshooting}

\subsection{Python Not Found}

Verify installation:

\begin{lstlisting}
python --version
\end{lstlisting}

or on macOS:

\begin{lstlisting}
python3 --version
\end{lstlisting}

\subsection{Git Not Found}

Verify installation:

\begin{lstlisting}
git --version
\end{lstlisting}

\subsection{Push Authentication Problems}

Use PyCharm's built-in GitHub login system rather than manually entering passwords or tokens.

\section{Final Verification Checklist}

Before class begins, verify that you can:

\begin{itemize}
    \item Run Python programs
    \item Open the repository in PyCharm
    \item Commit changes
    \item Push changes to GitHub
    \item Pull updates from GitHub
\end{itemize}

\end{document}
